\chapter{Постановка задачи}
\label{cha:Postanovka}

Задача автоматической паллетизации возникает в контексте складской логистики, где требуется разместить набор разнородных коробок на стандартизированном поддоне (паллете) ограниченных размеров. Качество полученной укладки определяется двумя конкурирующими критериями: плотностью заполнения объёма и физической устойчивостью штабеля. В данной работе решаются две связанные подзадачи.

\section{Подзадача 1. Укладка на одну паллету}

Дан набор коробок $S = \{b_1, b_2, \ldots, b_K\}$, каждая из которых характеризуется целочисленными размерами $(l_i, w_i, h_i)$ и массой $m_i$. Дана паллета с прямоугольным основанием $L \times W$ и нормировочной высотой $H$.

Требуется определить для каждой коробки координаты размещения $(x_i, y_i, z_i)$ и ориентацию таким образом, чтобы:
\begin{enumerate}
  \item все коробки помещались внутри области $[0, L] \times [0, W]$ по основанию;
  \item коробки не пересекались попарно;
  \item каждая коробка опиралась либо на поверхность паллеты ($z_i = 0$), либо на верхнюю грань хотя бы одной ранее уложенной коробки;
  \item коэффициент заполнения (percolation) был максимален;
  \item укладка удовлетворяла ограничению на устойчивость.
\end{enumerate}

Допустимыми являются только повороты вокруг вертикальной оси (перестановка $l_i$ и $w_i$), поскольку переворачивать коробки в практических сценариях недопустимо.

Ограничение устойчивости формализуется через минимальный коэффициент опоры (min\_support\_ratio). Для каждой коробки вычисляется доля площади её нижней грани, которая опирается на нижележащий слой. Минимальное значение этой доли по всем коробкам определяет устойчивость укладки в целом. Подробное обоснование выбора данной метрики приведено в разделе~\ref{sec:metric_support}.

\section{Подзадача 2. Укладка на несколько паллет}

В практической ситуации для одного адреса доставки может потребоваться несколько паллет. Дан тот же набор коробок $S$ и количество доступных паллет $P \geq 2$, каждая с одинаковым основанием $L \times W$.

Помимо требований, перечисленных для подзадачи~1, предъявляется дополнительное условие: высоты штабелей на разных паллетах должны быть по возможности одинаковыми. Неравномерная загрузка приводит к неэффективному использованию кузова транспортного средства и повышенному риску повреждения при перевозке.

Равномерность загрузки оценивается метрикой баланса высот (height\_balance), которая принимает значение~1 при одинаковых высотах и убывает при увеличении разброса (формальное определение --- в разделе~\ref{sec:metric_hb}).

\section{Связь с классическими задачами комбинаторной оптимизации}

Задача паллетизации является частным случаем задачи трёхмерной упаковки (3D Bin Packing Problem, 3D-BPP), которая принадлежит классу NP-трудных задач~\cite{Martello2000}. Типология Wäscher et al.~\cite{Wascher2007} относит данную постановку к классу \emph{Single Bin-Size Bin Packing Problem} при нескольких паллетах и \emph{Single Large Object Placement Problem} при одной паллете.

Существенным отличием от абстрактной 3D-BPP является наличие дополнительных практических ограничений: требование устойчивости, ограниченность поворотов, неоднородность коробок по весу и категории товара. Обзор таких ограничений систематизирован в работе~\cite{Bortfeldt2013}.
